% beamer-voorbeeld.tex
%
\documentclass{beamer} %OF \documentclass[handout]{beamer}
\usepackage[utf8]{inputenc}
\usepackage[dutch]{babel}
\usepackage{graphicx}
%%
%% Er zijn heel wat verschillende thema's ter beschikking!
%%
\usetheme{Madrid}
\usecolortheme{seahorse}
%
\setbeamercovered{transparent}
%% probeer ook eens invisible, transparent, dynamic of highly dynamic.
%% invisible is standaard
%
\title[Beamer]{Beamer: presentaties maken met \LaTeX}
\subtitle{Een instructief voorbeeld}
\date[2020--2021]{Computerproject Wiskunde}
\author[K. Coolsaet]{K. Coolsaet (G. Deschrijver)}
\institute[UGent]{Universiteit Gent}
%
\newtheorem{stelling}{Stelling}
\newtheorem{gevolg}[stelling]{Gevolg}
\newtheorem{voorbeeld}[stelling]{Voorbeeld}
\begin{document}

\begin{frame}[plain]
\titlepage
\end{frame}

%%%%%%%%%%%%%%%%%%%%%%%%%%%%%%%%%%%%%%%%%%%%%%%%%%%%%%%%%%%%%%%%%%%%%%
%\part{Introductie}
%\frame{\partpage}
%%%%%%%%%%%%%%%%%%%%%%%%%%%%%%%%%%%%%%%%%%%%%%%%%%%%%%%%%%%%%%%%%%%%%%

\section{Inleiding}

\begin{frame}{\secname} %% \secname bevat de laatste 'section'-titel
  \begin{itemize}
  \item Beamer laat toe om \emph{wiskundige} slides te maken, omdat je nog
  steeds alle standaard \LaTeX-opdrachten kunt gebruiken
  \pause
  \item Je kan ook `\emph{overlays}' maken.
  \pause
  \item Er bestaan verschillende Beamerstijlen om het
  uitzicht van de presentatie te veranderen.
  \pause
  \item Het output-bestand is een PDF-bestand dat op alle platformen exact
  hetzelfde uitzicht heeft.
  \end{itemize}
\end{frame}

%%%%%%%%%%%%%%%%%%%%%%%%%%%%%%%%%%%%%%%%%%%%%%%%%%%%%%%%%%%%%%%%%%%%%%

\section{Frames}

\begin{frame}{\secname}{Opties}

Een presentatie wordt opgedeeld in `\emph{frames}' (zie broncode).
  \pause  
  \begin{itemize}
  \item Frames kunnen een titel hebben en een subtitel, maar die zijn niet
  verplicht
  \pause
  \item Frames met optie [\texttt{plain}] tonen geen koplijnen, voetbalken,
  zijbalken, \dots
  \end{itemize}
\pause
Je kan nog steeds \texttt{section} en \texttt{subsection} gebruiken, en een
inhoudstafel laten maken, zoals in volgende slide.
\pause
Of in de slide daarna, waar we enkel de huidige paragraaf tonen
\end{frame}


\begin{frame}{Inhoud}
\tableofcontents  %% probeer ook eens \tableofcontents[pausesections]
\end{frame}

\begin{frame}{Inhoud --- tweede versie}
\tableofcontents[currentsection]
\end{frame}
%%%%%%%%%%%%%%%%%%%%%%%%%%%%%%%%%%%%%%%%%%%%%%%%%%%%%%%%%%%%%%%%%%%%%%

\section{Nieuwe omgevingen}
%
\begin{frame}{\secname}{Blokken}
  \begin{block}{Een blok}
  Dit blok werd aangemaakt met de \LaTeX-omgeving `\texttt{block}'.

  Een blok kan gerust meerdere paragrafen bevatten.
  \end{block}
  \pause
  \begin{alertblock}{Let op}
  En dit blok werd geproduceerd met `\texttt{alertblock}'. Hetzelfde
  kleureffect kan je ook \alert{`\emph{in line}'} verkrijgen.
  \end{alertblock}
  \pause
  \begin{exampleblock}{Voorbeeld}
  En met `\texttt{exampleblock}'
  \end{exampleblock}
  Titel is verplicht! (Maar maak hem desnoods blanco.)
\end{frame}

\begin{frame}{Stellingen}
Stellingen (en bewijzen) worden als blokken getoond.
\begin{stelling}[Murphy]
Alle oneven getallen zijn priem.
\end{stelling}
\begin{proof}
3 is priem, 5 is priem, 7 is priem, 9 is priem, 11 is priem, 13 is priem, enz.

De stelling volgt door inductie.
\end{proof}
\begin{voorbeeld}
113, 171 en 211 zijn priem.
\end{voorbeeld}
\begin{gevolg}
Geen enkel even getal is priem, behalve 2.
\end{gevolg}
\end{frame}


\section{Meerdere kolommen}

\begin{frame}{\secname}
  \begin{columns}
    \begin{column}[t]{5cm}
      \begin{block}{}
Tekst kan in meerdere kolommen worden opgedeeld.

Merk trouwens op dat dit blok in de linkerkolom \emph{geen} titel heeft.
      \end{block}
    \end{column}
    \pause
    \begin{column}[t]{5cm}
      In dit voorbeeld:
      \begin{itemize}
      \item 2 kolommen die 5cm breed zijn
      \item in de kolommen kan je tekst, blokken en andere omgevingen opnemen
      \end{itemize}
    \end{column}
  \end{columns}
\end{frame}

%%%%%%%%%%%%%%%%%%%%%%%%%%%%%%%%%%%%%%%%%%%%%%%%%%%%%%%%%%%%%%%%%%%%%%

\section{Overlays}

\begin{frame}{\secname}
Overlays kunnen nog ingewikkelder gemaakt worden:
\begin{itemize}
\item<4-> En nu pas zijn alle puntjes zichtbaar.
\item<3-> Toen kwam er een tweede lijn bij
\item<2-> Eerst zag je niets
\end{itemize}

\only<1-4>{Bijvoorbeeld wanneer je de dingen niet \textbf<2-4>{van boven naar
beneden} wil laten verschijnen.}
\begin{alertblock}<5->{Het kan trouwens nog veel erger?!}
Niet overdrijven!
\end{alertblock}
\end{frame}


\end{document}

